% Discussion body. Hand-written; numbers come from numbers.tex.

\noindent\textbf{What is proved, and what the measurements do not show.}
The Lean artifact certifies an implication: periodic conditioning plus
contraction implies a finite counting horizon. It does not prove any transformer
contracts, and two attempts to measure $q$ failed --- a boundary-distance
estimator needs a monotone read-head these decoders lack, and a perturbation
probe found no difference in local contraction between repeated and control items
($7$ of $18$ pairs, $p=0.12$), its own linearity control failing (S5).

The theorem's distinctive prediction is a constant $N^\ast$, invisible where our
ladder stopped; extending to $k=128$ found saturation of the predicted kind,
though not one derived from a measured $q$. A linear probe does not separate this
from the rival on which the count survives and only the output policy
fails~\cite{venkatesh2026repeated}: it retains less on repeated text in
\ProbeNBelowCtl{} of \ProbeNCk{} checkpoints, but the count stays decodable late
($R^2$ \ProbeLateMed{}), as that rival predicts.

Run past the horizon, the same probe does test the theorem, which forbids a
Lipschitz readout only \emph{beyond} $N^\ast$. On states from $k=48\ldots128$ it
reaches mean absolute error \PhPastMae{} in $\log_2 k$ against \PhPastConst{}
for a predictor that ignores the states and answers the mean --- learning nothing
--- where below the horizon it reaches \PhBelowMae{} against \PhBelowConst{}. We use that baseline because $R^2$ falls when the target varies
less, and the extension spans a third of the range. The theorem's
prediction, tested where it applies, on \PhPastN{} items and one checkpoint: one
clean observation, not a panel result (S12).
A conditional account risks being unfalsifiable, so we say what would refute it:
a control with its twin's deficit, a model whose capacity gain matches its
control, or a probe that recovers the count past the horizon.

The closest prior result is representational
collapse~\cite{barbero2024glasses}, where distinct prefixes converge as sequences
grow; our control separates that from periodicity. Related
bounds~\cite{dong2021attention,geshkovski2023emergence,gao2026clustered,%
strobl2024formal,merrill2024illusion} constrain one forward pass.

% Any further body text pushes the references onto a sixth page: main.tex ends
% with \vfill\pagebreak, so the bibliography always starts on a fresh page and
% a single spilled line costs a whole one. Check with build.sh before adding.
